\documentclass[../../../include/open-logic-section]{subfiles}

\begin{document}

\olfileid[fr]{sth}{replacement}{intro}
\olsection{Introduction}

Le remplacement est le schéma d'axiomes qui distingue $\ZF$ de~$\Z$.
Nous y avons eu recours tout au long des
\crefrange{sth:ordinals::chap}{sth:spine::chap}. Dans ce chapitre,
nous examinerons enfin la question suivante : le remplacement est-il justifié ?

Pour préciser cette question, remarquons que le remplacement est en réalité
assez \emph{puissant}. Ce chapitre nous donnera une idée de l'étendue de sa
puissance (de même que la \olref[sth][card-arithmetic][fix]{sec}). Mais cela
montrera aussi pourquoi il faut effectivement le justifier.

Nous examinerons deux sortes de justification. En gros, une justification
\emph{extrinsèque} cherche à justifier un axiome par ses fruits ; une
justification \emph{intrinsèque} cherche à le justifier en faisant valoir
qu'il se justifie à partir des concepts mathématiques en jeu. Ce chapitre nous permettra
de mieux comprendre cette distinction, qui n'est pourtant que la pointe
de l'iceberg. Pour aller plus loin, voir notamment \citeauthor{Maddy1988a}
(\citeyear{Maddy1988a} et \citeyear{Maddy1988b}).

\end{document}
